\section{Conclusiones}

El análisis integral de la incidencia delictiva en México entre 2015 y 2025 demuestra que la criminalidad
no es un fenómeno aleatorio, sino que sigue patrones sistemáticos determinados por la ubicación geográfica,
la clasificación delictiva y las condiciones temporales. El modelo predictivo de Random Forest alcanzó un R$^2$
de 0.93, lo que confirma que el 93\% de la variación en la incidencia delictiva puede explicarse mediante la
entidad federativa, el tipo de delito y variables como año, mes y trimestre. Esto descarta la hipótesis de
caos o impredecibilidad: los patrones delictivos son estables y recurrentes bajo condiciones normales, lo
cual es relevante para la planeación de políticas públicas de seguridad.

El análisis exploratorio reveló que la incidencia delictiva no ha disminuido en más de una década, con un
incremento notable a partir de 2022, y que los picos se concentran entre marzo y junio en los años recientes.
El robo es el delito más reportado a nivel nacional, mientras que estados como el Estado de México, CDMX y
Jalisco lideran en volumen absoluto. El clustering identificó cuatro perfiles delictivos claramente
diferenciados: un perfil de violencia contra las personas (sur y centro-pacífico), uno patrimonial
(áreas urbanas de alta densidad), uno familiar y de libertad personal, y uno diversificado. Estos perfiles
permiten entender que cada entidad federativa enfrenta dinámicas criminales distintas que requieren
estrategias de prevención diferenciadas.

Entre las limitaciones del proyecto destacan que el modelo aprende de patrones históricos y no puede
anticipar eventos anómalos como pandemias o crisis políticas, y que los datos dependen de denuncias
ante el Ministerio Público, lo que puede generar subreporte en ciertas zonas. Además, el modelo identifica
correlaciones, no causalidades: saber que CDMX tiene alta incidencia de robo no explica por qué ocurre.

No obstante, el proyecto logra su objetivo principal de proporcionar una herramienta analítica que permite
identificar los delitos más probables por estado y predecir su incidencia, ofreciendo una base sólida para
la toma de decisiones en materia de seguridad pública.